% Generated by scripts/assemble.py; corrections belong in source/corrections.json.
\examhead{22}{中国矿业大学 2024 年数学分析试卷}{26}
% p027-b07; source PDF page 27
\blockanchor{p027-b07}
\question{1}(15 分) 设定义在 $[a, b]$ 上的连续函数 $f(x)$ 满足 $f(a) = f(b)$, 且 $f_{+}^{\prime}(a)f_{-}^{\prime}(b) > 0$. 证明: 在 $(a, b)$ 内至少有一点 $\xi$, 使得 $f(\xi) = f(a)$.\par
% p027-b08; source PDF page 27
\blockanchor{p027-b08}
\question{2}(15 分) 叙述并证明 $\lim_{x \to +\infty} f(x) = +\infty$ 的归结原则.\par
% p027-b09; source PDF page 27
\blockanchor{p027-b09}
\question{3}(15 分) 叙述有限覆盖定理和致密性定理, 并用有限覆盖定理证明致密性定理.\par
% p027-b10; source PDF page 27
\blockanchor{p027-b10}
\question{4}(15 分) 求极限 $\lim_{x \to 0} \frac{1}{x} \left( \frac{1}{x} - \cot x \right)$.\par
% p027-b11; source PDF page 27
\blockanchor{p027-b11}
\question{5}(15 分) 设 $f(x)$ 与 $g(x)$ 均是定义在 $[a, b]$ 上的可积函数, 证明: 函数 $f(x)g(x)$ 也在 $[a, b]$ 上可积.\par
% p027-b12; source PDF page 27
\blockanchor{p027-b12}
\question{6}(15 分) 求 $I = \int_{0}^{1} tg(t) \, \mathrm{d}t$, 其中 $g(t) = \int_{t^2}^{1} e^{-x^2} \, \mathrm{d}x$.\par
% p027-b13; source PDF page 27
\blockanchor{p027-b13}
\question{7}(15 分) 求符号函数 $f(x) = \operatorname{sgn} x$ 在 $(- \pi, \pi)$ 上的傅里叶级数并判断其敛散性.\par
% p027-b14; source PDF page 27
\blockanchor{p027-b14}
\question{8}(15 分) 证明当 $\alpha > 0$ 时, 函数项级数 $\sum_{n=1}^{\infty} \frac{\sin nx}{n^{2+\alpha}}$ 在 $(- \infty, +\infty)$ 上一致收敛, 且其和函数 $S(x)$ 在 $(- \infty, +\infty)$ 内具有连续的导函数.\par
% p027-b15; source PDF page 27
\blockanchor{p027-b15}
\question{9}(15 分) 设二元函数组 $\left\{ \begin{array}{l} u = f(x, y); \\ v = g(x, y). \end{array} \right.$ 将 $xOy$ 平面上的点集 $E$ 映射成 $uOv$ 平面上的点集 $D$. $f(x, y)$ 和 $g(x, y)$ 为点集 $E$ 上的一致连续函数, 且 $h(u, v)$ 为点集 $D$ 上的一致连续函数, 证明: 复合函数 $h(f(x, y), g(x, y))$ 为点集 $E$ 上的一致连续函数.\par
% p027-b16; source PDF page 27
\blockanchor{p027-b16}
\question{10}(15 分) 设 $I(\alpha) = \int_{0}^{+\infty} e^{-\alpha x} \frac{\sin \beta x}{x} \, \mathrm{d}x$, 其中 $\beta > 0$ 为常数. 证明:\par
% p027-b17; source PDF page 27
\blockanchor{p027-b17}
(1) (7 分) $I(\alpha)$ 在 $(0, +\infty)$ 上连续.\par
% p027-b18; source PDF page 27
\blockanchor{p027-b18}
(2) (5 分) $I(\alpha)$ 在 $(0, +\infty)$ 上可导并求 $I(\alpha)$ 的表达式.\par
% p027-b19; source PDF page 27
\blockanchor{p027-b19}
(3) (3 分) 利用 $I(\alpha)$ 的表达式求 $D(\beta) = \int_{0}^{+\infty} \frac{\sin \beta x}{x} \, \mathrm{d}x$ 的值.\par
